\documentclass[10pt]{article}
\usepackage[margin=20mm]{geometry}
\usepackage{amsmath,amssymb,mathtools,fontspec,unicode-math,xurl,hyperref}
\IfFontExistsTF{DejaVu Serif}{\setmainfont{DejaVu Serif}}{\setmainfont{Latin Modern Roman}}
\setmathfont{latinmodern-math.otf}
\setlength{\parindent}{0pt}\setlength{\parskip}{6pt}\sloppy
\title{Technical Review of Distances and Isoperimetric Inequalities in Random Triangulations of High Genus}
\author{AI-assisted technical review}
\date{}
\begin{document}
\maketitle
Abstract: This review concerns arXiv:2311.04005v1. The exponential estimate in Lemma 14 uses the larger of two step probabilities as a lower-bound base, and its displayed inequality cannot hold with a fixed positive constant. This step supports the upper Cheeger bound and the separation between diameter and typical distance. A uniform corrected estimate restores the required logarithmic scale, conditional on the preceding encoding and conditioning arguments. A separate domain error in a distance decomposition is confined by the exceptional-event restriction used later in the proof. Neither finding refutes the final theorems.\par
\section*{Identity and inspected scope}
The paper by Thomas Budzinski, Guillaume Chapuy, and Baptiste Louf is associated with The Annals of Probability 53 (2025), 1645–1667, DOI 10.1214/24-AOP1746. The reviewed artifact is the 23-page author preprint [paper], arXiv:2311.04005v1, submitted on 7 November 2023; its title-page body date is 8 November 2023. Page numbers below are one-based PDF pages, equal to the printed pages. Equivalence with the journal text is not asserted.\par
The full extracted main text and references were read. The detailed audit concentrates on Lemma 14’s probability calculation and its use in Theorem 4 and Proposition 17, and on the domain of equation (16). Definitions and Theorems 1–4 were checked; the enumerative, isoperimetric, and ball-growth arguments were followed for dependency and scope. External enumerative results and the full tree-encoding literature were not independently re-proved. This is a claim-focused initial review.\par
\section*{Definitions and operative assumptions}
The model is a uniform rooted triangulation \(T_{2n,g_n}\) of an oriented surface, with \(2n\) faces, loops and multiple edges allowed, and \(g_n/n\to\theta\in(0,1/2)\) [Retained latex: \textbackslash{}cite[pp.\textasciitilde{}2--4]\{paper\}]. Tentacles arise by recursive insertions into a digon; the core is obtained by removing maximal tentacles [Retained latex: \textbackslash{}cite[pp.\textasciitilde{}16--17]\{paper\}]. Write \(T_n=an\) and \(M_n=2bn\) for their number and total number of internal faces. Lemma 13 supplies positive lower bounds for \(a,b\) with high probability. The deterministic bounds \(a\leq3\), \(b\leq1\) suffice below, since the original map has \(3n\) edges and \(2n\) faces.\par
Equation reference: eq:pq. Conditionally on these counts and the core, Lemma 14 uses a walk of length \(L=(a+3b)n\). The independent comparison walk \(W\) has increments \(+2,-1\) with respective probabilities \[
p=\frac b{a+3b},\qquad q=\frac{a+2b}{a+3b}.\] Every admissible positive \(a,b\) satisfies \(p+q=1\) and \(0<p<q<1\). Count integrality is inherited from the conditioning; the analytic argument does not posit a nonintegral triangulation. The conditioned walk \(V\) is obtained by an event with probability \(\Omega(n^{-1/2})\), according to the endpoint estimate and cycle-lemma calculation on page 18.\par
Equation reference: eq:u. For odd \(\ell\), the prescribed word \(w_\ell\) contains \(\ell-1\) up-steps and \(2\ell-1\) down-steps. Its length and probability are therefore \[
|w_\ell|=3\ell-2,\qquad u_\ell=p^{\ell-1}q^{2\ell-1}.\] These are the source’s pattern and hypotheses.\par
\section*{Substantive findings}
Problem 1 (Incorrect exponential base in equation (12)).  The second inequality of equation (12), in Lemma 14 on page 18, is false with the printed \(c=q=(a+2b)/(a+3b)\) and a fixed \(c'>0\).\par
\subsection*{Original claim and decisive proof}
Equation reference: eq:required, eq:ratio. Disjoint blocks give the valid estimate \[\mathbb P(W\hbox{ avoids }w_\ell)\leq(1-u_\ell)^{N_\ell},
\qquad N_\ell=\left\lfloor\frac{(a+3b)n}{3\ell-2}\right\rfloor.\] The source then claims that this last expression is at most \(\exp\{-c'nq^{3\ell}/\ell\}\). Taking negative logarithms would require \[
\frac{N_\ell[-\log(1-u_\ell)]}{(n/\ell)q^{3\ell}}\geq c'.\] For fixed admissible positive \(a,b\) and odd \(\ell\to\infty\) with \(\ell=o(n)\), the left-hand side tends to zero: \[
\frac{u_\ell}{q^{3\ell}}
=\frac1{pq}\left(\frac pq\right)^\ell\longrightarrow0,\quad
\frac{\ell N_\ell}{n}\longrightarrow\frac{a+3b}{3},\quad
\frac{-\log(1-u_\ell)}{u_\ell}\longrightarrow1.\] No fixed positive \(c'\) repairs this inequality with the printed base. Uniformly on the positive lower-count event, \(p/q\leq1/2\) and \(1/(pq)\) is bounded, so the same contradiction holds when \(a,b\) vary with \(n\). The contradiction is algebraic throughout the positive parameter domain, not a finite simulation or a selected unusual triangulation.\par
The claimed condition on \(s_\theta\) is also insufficient. Because \(pq^2<q^3\), there exist \[\frac1{\log(1/(pq^2))}<s<\frac1{3\log(1/q)}.\] For odd \(\ell=s\log n+O(1)\), a union bound over all starting positions gives \[\mathbb P(W\hbox{ contains }w_\ell)\leq Lu_\ell
=O\!\left(n^{1-s\log(1/(pq^2))}\right)\longrightarrow0.\] Avoidance thus tends to one for this independent walk, whereas the claimed upper bound tends to zero. This refutes the intermediate estimate, not the existence assertion of Lemma 14.\par
\subsection*{Independent computation and counterevidence}
Exact-rational arithmetic for the illustrative walk parameters \(a=3/2,b=1/2\), hence \(p=1/6,q=5/6\), yields \[u_5/q^{15}=36/15625=0.002304,\qquad
u_{11}/q^{33}=36/244140625=1.47456\times10^{-7}.\] The two thresholds above are approximately \(0.4637353\) and \(1.8282716\). The retained calculation is corroborative; these exact counts are not asserted to have been sampled from a triangulation.\par
The actual page image confirms that \(c=q\) is printed, excluding an extraction reversal. Permitting \(c'\) to vanish with \(\ell\) removes the constant needed in the proof. Interpreting \(c\) as the smaller probability \(p\) does yield a valid local repair, so a wrong choice of constant is a plausible explanation. It does not validate the displayed inequality.\par
\subsection*{Uniform repair}
Equation reference: eq:repair. On the high-probability count event, take \(a\geq a_0>0\), \(b\geq b_0>0\). Then \[p\geq p_0:=b_0/6>0,\qquad q\geq p\geq p_0.\] Decrease \(p_0\) below \(1\) if necessary. For \(\ell=o(n)\), uniformly for sufficiently large \(n\), \[N_\ell\geq\kappa n/\ell,\quad \kappa=(a_0+3b_0)/6,\qquad
u_\ell\geq p_0^{3\ell-2}\geq p_0^{3\ell}.\] Consequently \[
\mathbb P(W\hbox{ avoids }w_\ell)
\leq\exp\!\left\{-\kappa(n/\ell)p_0^{3\ell}\right\}.\] Choose \(0<s<1/[3\log(1/p_0)]\) and odd \(\ell=s\log n+O(1)\). The exponent magnitude is at least a positive multiple of \(n^\varepsilon/\log n\) for some \(\varepsilon>0\). Dividing this failure probability by the conditioning probability \(\Omega(n^{-1/2})\) still gives zero in the limit. This supplies the missing estimate for a logarithmic tentacle, assuming the paper’s preceding encoding and conditioning construction. It preserves existence of a positive constant, without establishing the larger printed range or independently proving every combinatorial input.\par
\subsection*{Dependency, severity, and surviving results}
Lemma 14 supplies the upper bound \(h(T_{2n,g_n})=O(1/\log n)\) in Theorem 4 (page 18) and \(H(\vec e^{\,*})\geq(s_\theta/2)\log n\) in Proposition 17 (page 20). The assessment is Major: invalid core-related derivation. A short repair preserves the qualitative conclusions; this bounds the impact without making the original derivation valid. No counterexample to the final graph theorems is established. The isoperimetric theorem and logarithmic diameter bounds precede this calculation and do not depend on it.\par
Problem 2 (The distance decomposition needs distinct tentacles).  Equation (16), introduced for arbitrary two oriented edges on page 20, cannot hold for two edges within the same deep tentacle.\par
\subsection*{Assumptions and complete local counterexample}
The displayed distance is the sum of two heights, the core distance between their bases, and an error in \([-2,2]\). Take an oriented edge \(\vec e\) with \(H(\vec e)>1\) inside a sufficiently long ladder, and put \(\vec e_1=\vec e_2=\vec e\). Both the original distance and core distance are zero. The equation requires error \(-2H(\vec e)<-2\), contrary to its bound. Recursive ladder insertion is allowed and preserves genus. More explicitly, fix a sufficiently deep finite ladder with \(2k\) faces and choose a genus-\(g_n\) core with \(2(n-k)\) faces. The necessary condition \(n-k\geq2g_n-1\) holds eventually because \(g_n/n\to\theta<1/2\); insertion produces the required \(2n\)-face triangulation. If the phrase two edges is interpreted as distinct, use two different outgoing oriented edges at the same sufficiently deep interior ladder vertex. Their starting-point distance is zero, their heights are large, and the core distance is nonnegative, again contradicting the bound.\par
\subsection*{Alternative reading, correction, and bounded impact}
For different tentacles, paths must leave each through its boundary endpoints, which is the intended setting of the core-plus-heights estimate. Page 21 explicitly discards the event that two selected edges belong to the same tentacle. This is counterevidence against treating the overstatement as a failure of Proposition 17. Restrict equation (16) to distinct tentacles and impose that restriction before its uses. The assessment is Minor: local domain overstatement. No additional central damage is demonstrated; the correction narrows the domain, without refuting the final proposition.\par
\section*{Version checks and conclusion}
A bounded official/author-source search did not identify a separate formal erratum. The current arXiv abstract fetch failed, so the complete current revision history was not independently recovered. An opened author-hosted PDF is also dated 8 November 2023 and has 23 pages. These observations do not prove absence of corrections. Journal-final proof text was not compared; the findings apply to the acquired version.\par
Item | Assessment | Repair and impact boundary\par
Problem 1 | Major | Corrected uniform exponential bound proved; final theorem not refuted.\par
Problem 2 | Minor | Restrict to distinct tentacles; later exceptional-event treatment supports intended use.\par
Decision | Pending | Initial AI review; independent acceptance not performed.\par
Adoption | Unknown | No author contact or adoption requested.\par
\section*{Supporting checks}
`e5e9ffb2f4b56fdfd3d2410fa7bba5da3557af8a3a5816f74bade943596ced6e`.\par
Source manifests, exact report/build bindings, compiler dependencies, and per-page render/view records are retained under this paper’s numeric paths. `evidence\_problem\_01.py` and `problem\_01\_evidence.json` contain the independent rational arithmetic; `problem\_02\_evidence.json` records the analytic counterexample. No author code or simulation was run.\par
The primary scientific executor read the paper and carried out the proof audit. Native helpers checked source bytes/hashes and discovered revision metadata only. Controller UI observations and the executor’s runtime self-report are separately retained in provenance, without independent backend certification. Human reviews remain empty and scientific adjudication remains pending.\par
9 T. Budzinski, G. Chapuy, and B. Louf, Distances and isoperimetric inequalities in random triangulations of high genus, arXiv:2311.04005v1 (7 November 2023). https://arxiv.org/abs/2311.04005v1. Journal DOI: https://doi.org/10.1214/24-AOP1746. T. Budzinski, author-hosted PDF of the same title, dated 8 November 2023, accessed 29 October 2026 for version metadata. https://perso.ens-lyon.fr/thomas.budzinski/papers/diameter\_high\_genus.pdf.\par
\end{document}
